\section{Push of the past and Pull of the present}


With regards to the standard birth death process in which the condition of survival to the present is imposed (Quoting R.  Mann and G. Budd, who themselves were overviewing the topic introduced by Nee et.  al.):


\begin{quotation}
\noindent
"The POTPa emerges as a feature of diversification by the fact
that all modern clades (tautologically) survived until the present
day. This singles them out from the total population of clades
that could be generated from any particular pair of background
birth and death rates: clades that happened by chance to start off
with higher net rates of diversification have better-than-average
chances of surviving until the present day. As Nee et al. (1994a)
put it, such clades “got off to a flying start”; and they accumu-
late species faster than one would expect. Once clades become
established, they are less vulnerable to random changes in the
observed diversification rate, and this value therefore tends to re-
vert to the background rate through time. Long-lived clades thus
tend to show high observed rates of diversification at their origin,
which then decrease to their long-term average as the present is
approached."
\end{quotation}

\noindent
In their paper, "History is written by the victors," R. Mann and G. Budd propose that the push of the past (POTPa) is an underappreciated phenomenon in the paleontological literature, which delves into mass radiations following mass extinctions.
Without directly opposing their work, I present a model that illustrates a dynamic of variable-rate evolution, which also seeks to shed light on the higher origination rates of species following extinctions and events of ecological instability more generally \cite{bak1991self} \cite{fisher1929genetical}\cite{budd2018history}.
